%=====================================================================
\section{Lecture 3: The Keldysh contour is already in the von Neumann equation}
\label{sec:lec3}

The Keldysh formalism is usually introduced by declaring a closed time contour and then building Green's functions on it.
In this lecture we go the other way: we expand the reduced density matrix in powers of $H_I$ and watch the contour, the Keldysh rotation and the $2\times2$ matrix structure of the bath propagator appear by themselves.
Nothing is postulated.
At the end we have a set of rules (``contractions'') that \cref{sec:lec4} turns into diagrams.

\begin{figure}[t]
\centering
\begin{tikzpicture}[>=Stealth,scale=1.0]
  \draw[thick,RoyalBlue,postaction={decorate},decoration={markings,mark=at position 0.5 with {\arrow{>}}}]
      (0,0.35) -- (6,0.35);
  \draw[thick,RoyalBlue] (6,0.35) arc (90:-90:0.35);
  \draw[thick,RoyalBlue,postaction={decorate},decoration={markings,mark=at position 0.5 with {\arrow{>}}}]
      (6,-0.35) -- (0,-0.35);
  \node[left] at (0,0.35) {$+$ (ket, $U$)};
  \node[left] at (0,-0.35) {$-$ (bra, $U^\dagger$)};
  \draw[->,gray] (0,-1.0) -- (6.8,-1.0) node[right] {$t$};
  \node[below,gray] at (0,-1.0) {$0$};
  \node[below,gray] at (6,-1.0) {$t$};
  \node[right,align=left,font=\small] at (7.0,0.0) {$\rho(t)=U(t)\rho(0)U^\dagger(t)$:\\
  ordering superoperators $A^\pm(t_i)$\\ on one real axis $\Leftrightarrow$ contour ordering};
\end{tikzpicture}
\caption{The Keldysh contour~\cite{Schwinger1961,Keldysh1964}.
The forward branch carries the operators of $U(t)$ acting on the ket side of $\rho(0)$, the backward branch the operators of $U^\dagger(t)$ acting on the bra side.
Contour ordering means: later contour position further from $\rho(0)$.}
\label{fig:contour}
\end{figure}

\subsection{The Dyson series for \texorpdfstring{$\rho$}{rho} and the origin of the contour}
\label{sec:contour-from-vN}

Work in the interaction picture with respect to $H_0=H_S+H_B$, \cref{sec:pictures}: $\tilde\rho(t)=\ee^{iH_0t}\rho(t)\ee^{-iH_0t}$ obeys $\dot{\tilde\rho}=-i[\tilde H_I(t),\tilde\rho]$ with $\tilde H_I(t)=\ee^{iH_0t}H_I\ee^{-iH_0t}$, \cref{eq:int-picture-vN}.
Integrating, $\tilde\rho(t)=\tilde U(t)\rho(0)\tilde U^\dagger(t)$ with
\begin{equation}
  \tilde U(t)=T\exp\Big[-i\!\int_0^t\!\dd s\,\tilde H_I(s)\Big]
  =1-i\!\int_0^t\!\dd t_1\tilde H_I(t_1)-\!\int_0^t\!\dd t_1\!\int_0^{t_1}\!\dd t_2\,\tilde H_I(t_1)\tilde H_I(t_2)+\dots
\end{equation}
Its adjoint is \emph{anti}-time-ordered: in $\tilde U^\dagger$ the later operator stands to the right, $\tilde U^\dagger=1+i\int\tilde H_I-\int_{t_1>t_2}\tilde H_I(t_2)\tilde H_I(t_1)+\dots$

\begin{exbox}[Example 3.1: the second-order term, written out]
Collecting all terms of order $H_I^2$ in $\tilde U\rho(0)\tilde U^\dagger$,
\begin{align}
  \tilde\rho^{(2)}(t)=&-\int_0^t\!\dd t_1\!\int_0^{t_1}\!\dd t_2\,\tilde H_I(t_1)\tilde H_I(t_2)\,\rho(0)
   &&\text{(both on the ket side; time-ordered)}\notag\\
  &-\int_0^t\!\dd t_1\!\int_0^{t_1}\!\dd t_2\,\rho(0)\,\tilde H_I(t_2)\tilde H_I(t_1)
   &&\text{(both on the bra side; anti-time-ordered)}\notag\\
  &+\int_0^t\!\dd t_1\!\int_0^{t}\!\dd t_2\,\tilde H_I(t_1)\,\rho(0)\,\tilde H_I(t_2)
   &&\text{(one on each side; \emph{no} mutual ordering)}.
  \label{eq:rho2}
\end{align}
The third term is the one that looks unusual: $t_1$ and $t_2$ are integrated independently over $[0,t]$, so there is no time ordering between a ket-side and a bra-side operator.
All three terms become a \emph{single} ordering rule once we say that every point on the forward ($+$) branch is ``earlier'' than every point on the backward ($-$) branch, that time runs forward on $+$ and backward on $-$, and that operators are ordered with later contour positions further from $\rho(0)$.
That rule is contour ordering, \cref{fig:contour}.
\end{exbox}

The superoperator language of \cref{sec:lec1} makes the rule automatic.
With $A^{+}X=AX$ and $A^{-}X=XA$,
\begin{equation}
  \tilde\rho^{(2)}(t)=\int_0^t\!\dd t_1\!\int_0^{t_1}\!\dd t_2\,
  \tilde\Lc_I(t_1)\tilde\Lc_I(t_2)\,\rho(0),\qquad
  \tilde\Lc_I(s)=-i\big[\tilde H_I^{+}(s)-\tilde H_I^{-}(s)\big],
  \label{eq:rho2super}
\end{equation}
and multiplying out the two brackets reproduces the three terms of \cref{eq:rho2}: $H^+H^+$ gives the first, $H^-H^-$ gives the second (note $H^-(t_1)H^-(t_2)\rho=\rho H(t_2)H(t_1)$, automatically anti-ordered), and the cross terms $H^+(t_1)H^-(t_2)+H^-(t_1)H^+(t_2)$ give $H(t_1)\rho H(t_2)+H(t_2)\rho H(t_1)$ restricted to $t_1>t_2$, which is the third term with its two orderings separated.
To all orders,
\begin{equation}
  \tilde\rho(t)=T\exp\Big[\int_0^t\!\dd s\,\tilde\Lc_I(s)\Big]\rho(0)
  =\sum_n\int_{t>t_1>\dots>t_n>0}\tilde\Lc_I(t_1)\cdots\tilde\Lc_I(t_n)\,\rho(0),
  \label{eq:dyson-super}
\end{equation}
which is just the Dyson series for the superoperator equation $\dot{\tilde\rho}=\tilde\Lc_I(t)\tilde\rho$ of \cref{eq:int-picture-super}, with $\tilde\Lc_I(t)=\ee^{-\Lc_0t}\Lc_I\ee^{\Lc_0t}$; that this agrees with the operator definition $\tilde\Lc_I=-i[\tilde H_I^+-\tilde H_I^-]$ used in \cref{eq:rho2super} is Lemma~\ref{lem:int-picture}(ii).

\begin{keybox}[Key point 3.1]
Ordinary time ordering of the superoperators $\tilde\Lc_I(t_i)$ along the single real axis $[0,t]$, with a sum over the branch index $\pm$ at each vertex, \emph{is} Keldysh contour ordering.
One never has to think about the contour again: it is encoded in the sign of $\Lc_I$ and the left/right action of $H^\pm$.
Conversely, every Keldysh expression for a one-time expectation value can be read as a Liouville-space string.
\end{keybox}

\subsection{The Keldysh rotation in Liouville space}
\label{sec:rotation-liouville}

The $\pm$ basis is natural for deriving the contour but awkward for calculations, because the two branches are strongly correlated (the trace at the final time glues them together).
The rotation \cref{eq:clqu} separates ``what the two branches have in common'' ($\cl$) from ``how they differ'' ($\qu$).
Since $S$ and $B$ act on different tensor factors, $(SB)^\pm=S^\pm B^\pm$ and $S^\pm$ commutes with $B^\pm$.
Then with $A^\pm=A^\cl\pm\tfrac12A^\qu$,
\begin{align}
  S^{+}B^{+}-S^{-}B^{-}
  &=\big(S^\cl+\tfrac12S^\qu\big)\big(B^\cl+\tfrac12B^\qu\big)
   -\big(S^\cl-\tfrac12S^\qu\big)\big(B^\cl-\tfrac12B^\qu\big)
  =S^\cl B^\qu+S^\qu B^\cl ,
\end{align}
the $S^\cl B^\cl$ and $\frac14S^\qu B^\qu$ terms cancelling.
Hence
\begin{equation}
  \boxed{\;\Lc_I=-i\big(S^{\qu}B^{\cl}+S^{\cl}B^{\qu}\big).\;}
  \label{eq:LIcq}
\end{equation}
The structure ``one $\cl$, one $\qu$ per vertex'' is the same as in the Keldysh action, where the interaction term $H_+-H_-$ is odd in the quantum field and so always contains at least one $\qu$ leg~\cite[Sec.~2.7]{Kamenev2011}.

\paragraph{The two trace identities.}
For an arbitrary operator $Y$ on $\mathcal H_S\otimes\mathcal H_B$ and a bath operator $B$,
\begin{equation}
  \Tr_B\!\big(B^\qu Y\big)=\Tr_B\!\big(BY-YB\big)=0,\qquad
  \Tr_B\!\big(B^\cl Y\big)=\tfrac12\Tr_B\!\big(BY+YB\big)=\Tr_B\!\big(BY\big),
  \label{eq:trace-identities}
\end{equation}
both by cyclicity of the partial trace over the bath factor.
The first identity is the superoperator form of the Keldysh statement that a correlation function whose \emph{latest} point is on a quantum leg vanishes; it is what makes $\avg{\phi_\qu\phi_\qu}=0$ and $\avg{\phi_\cl\phi_\qu}$ retarded~\cite[Sec.~2.6]{Kamenev2011}.
The same identities hold with $S$ in place of $B$ and $\Tr_S$ in place of $\Tr_B$, and they will be used on both sides.

\begin{exbox}[Example 3.2: what $\cl$ and $\qu$ do to a qubit state]
From \cref{eq:SclSqu}, acting on $\rho=\tfrac12(\Id+\bm r\cdot\bm\sigma)$:
\begin{equation}
  S^\cl\rho=\tfrac12\big[(\hat n\cdot\bm r)\Id+\hat n\cdot\bm\sigma\big],\qquad
  S^\qu\rho=i\,(\hat n\times\bm r)\cdot\bm\sigma,\qquad \hat n=(\sin\theta,0,\cos\theta).
\end{equation}
$S^\cl$ produces a traceful object whose weight is the expectation value $\avg{S}=\hat n\cdot\bm r$; $S^\qu$ produces a traceless object proportional to the torque $\hat n\times\bm r$ that a field along $\hat n$ would exert.
A state with $\bm r\parallel\hat n$ (a pointer state of the coupling) has $S^\qu\rho=0$: it does not ``differ between the branches''.
For $\theta=0$ this is the statement that a diagonal density matrix has no quantum component with respect to $\sz$.
This is the operator version of ``the classical field is the physical configuration, the quantum field measures the difference between the two branches''.
\end{exbox}

\subsection{Bath contractions and the \texorpdfstring{$2\times2$}{2x2} Keldysh structure}
\label{sec:contractions}

The basic object is the time-ordered two-point function of bath superoperators,
\begin{equation}
  D^{ab}(t,s)\equiv\Tr_B\!\big[T_{\!\Lc}\,\tilde B^{a}(t)\tilde B^{b}(s)\,\rhoB\big],\qquad a,b\in\{+,-\}\ \text{or}\ \{\cl,\qu\},
\end{equation}
where $T_{\!\Lc}$ puts the later superoperator on the left (closer to the trace).
In the $\pm$ basis the four components are, by direct evaluation for $t>s$,
\begin{center}
\renewcommand{\arraystretch}{1.3}
\begin{tabular}{@{}llll@{}}
\toprule
component & string ($t>s$) & value & Keldysh name \\
\midrule
$D^{++}(t,s)$ & $\Tr_B[B(t)B(s)\rhoB]$ & $C(t-s)$ & time-ordered $D^T$\\
$D^{--}(t,s)$ & $\Tr_B[\rhoB B(s)B(t)]$ & $C(s-t)$ & anti-time-ordered $D^{\tilde T}$\\
$D^{+-}(t,s)$ & $\Tr_B[B(t)\rhoB B(s)]$ & $C(s-t)$ & lesser $D^{<}$\\
$D^{-+}(t,s)$ & $\Tr_B[B(s)\rhoB B(t)]$ & $C(t-s)$ & greater $D^{>}$\\
\bottomrule
\end{tabular}
\end{center}
(For $t<s$ the roles of ``later'' and ``earlier'' are exchanged in the string; the $+-$ and $-+$ entries then keep their form while $++$ and $--$ switch to $C(s-t)$ and $C(t-s)$ respectively.)
These are the four contour-ordered Green's functions of Keldysh theory up to the overall $-i$ of the usual definition, and the familiar linear relation $D^T+D^{\tilde T}=D^<+D^>$ holds entry by entry, which is why the $\pm$ basis is redundant.
Now rotate, using $B^\cl=\frac12(B^++B^-)$ and $B^\qu=B^+-B^-$ on both legs:
\begin{align}
  D^{\cl\cl}(t,s)&=\tfrac14\big(D^{++}+D^{+-}+D^{-+}+D^{--}\big)
   =\tfrac12\big[C(t-s)+C(s-t)\big]=S(t-s),\\
  D^{\cl\qu}(t,s)&=\tfrac12\big(D^{++}-D^{+-}+D^{-+}-D^{--}\big)
   =\theta(t-s)\avg{[B(t),B(s)]}=-i\chi(t-s),\\
  D^{\qu\cl}(t,s)&=\tfrac12\big(D^{++}+D^{+-}-D^{-+}-D^{--}\big)
   =-\theta(s-t)\avg{[B(t),B(s)]}=-i\chi(s-t),\\
  D^{\qu\qu}(t,s)&=D^{++}-D^{+-}-D^{-+}+D^{--}=0 .
\end{align}
(For the middle two lines evaluate separately for $t>s$ and $t<s$ using the table; the $\theta$ functions appear because the ``later'' entry of the table switches.)
In matrix form,
\begin{equation}
  \boxed{\;
  D(t,s)=\begin{pmatrix} D^{\cl\cl} & D^{\cl\qu}\\ D^{\qu\cl} & D^{\qu\qu}\end{pmatrix}
   =\begin{pmatrix} S(t-s) & -i\chi(t-s)\\ -i\chi(s-t) & 0\end{pmatrix}
   \;\leftrightarrow\;
   \begin{pmatrix} D^K & D^R\\ D^A & 0\end{pmatrix},\;}
  \label{eq:LO}
\end{equation}
the Larkin--Ovchinnikov structure~\cite{LarkinOvchinnikov1975,RammerSmith1986}. The zero
in the $\qu\qu$ slot and the retardation of the $\cl\qu$ slot are both consequences of the
first trace identity in \cref{eq:trace-identities}: whenever the \emph{later} superoperator
in the string is $B^\qu$, the trace kills the term. Restricting to $t>s$ and reading off the
entries gives the three contraction rules used from now on:
\begin{equation}
  \Tr_B[\tilde B^{\cl}(t)\tilde B^{\cl}(s)\rhoB]=S(t-s),\qquad
  \Tr_B[\tilde B^{\cl}(t)\tilde B^{\qu}(s)\rhoB]=-i\chi(t-s),\qquad
  \Tr_B[\tilde B^{\qu}(t)\,\cdot\,]=0 .
  \label{eq:contractions}
\end{equation}

\begin{exbox}[Example 3.3: the single-mode bath, and the Lorentzian]
For one mode, $C(t)=g^2[(\bar n+1)\ee^{-i\omega_bt}+\bar n\,\ee^{i\omega_bt}]$, the table gives
$D^{\cl\cl}=g^2(2\bar n+1)\cos\omega_b(t-s)$ and $D^{\cl\qu}=-2ig^2\theta(t-s)\sin\omega_b(t-s)$: the noise
line carries the occupation, the response line does not. In frequency space
$\chi''(\omega)=\pi g^2[\delta(\omega-\omega_b)-\delta(\omega+\omega_b)]$ and $S(\omega)=\coth(\beta\omega/2)\chi''(\omega)$, the FDT checked
by hand. For the Lorentzian bath, \cref{eq:lorentz-kernels} gives the same with
$\omega_b\to\Omega$, $\bar n\to0$ and an extra damping $\ee^{-\gamma|t-s|/2}$ on both lines. In
either case the Keldysh matrix \cref{eq:LO} is, in frequency space,
\begin{equation}
  D(\omega)=\begin{pmatrix} \coth\tfrac{\beta\omega}2\,\chi''(\omega) & -i\chi(\omega)\\ -i\chi(\omega)^* & 0\end{pmatrix},
\end{equation}
which is Kamenev's $D^K=\coth(\beta\omega/2)(D^R-D^A)$ with $D^R=-i\chi$~\cite[Sec.~2.6]{Kamenev2011}.
\end{exbox}

\paragraph{Wick's theorem for superoperators.}
Because the Liouville string $\Tr_B[\tilde B^{a_n}(t_n)\cdots\tilde B^{a_1}(t_1)\rhoB]$ with
$t_n>\dots>t_1$ is literally a contour-ordered expectation value (Key point 3.1), the Wick
theorem \eqref{eq:wick} holds for it unchanged, with each pair contributing
$D^{a_ia_j}(t_i,t_j)$ from \cref{eq:LO}. In particular, every pairing in which the latest
superoperator of a pair is $B^\qu$ vanishes.

\begin{warnbox}
\small\textbf{Two things that are easy to get wrong.} (1) The $\theta$ function in $\chi$ is
not an approximation; it comes from the trace identity and is the reason causality is
built in. If you ever find a $\chi(t-s)$ with $s>t$ in a system-side diagram, an index
has been swapped. (2) The $-i$ in the response line is a convention tied to
$\chi=i\theta\avg{[B,B]}$. With Kamenev's $D^R=-i\theta\avg{[\phi,\phi]}$ the line is
$+D^R$. Keep one convention and check it on the single-mode example above.
\end{warnbox}

\subsection{Exercises}

\begin{exercise}\label{ex:3-1}
Write out the third-order term of \cref{eq:dyson-super} in the $\pm$ basis and identify
the eight terms with the eight assignments of branches to three contour points. Show that
after $\Tr_B$ with a Gaussian bath all of them vanish (odd moments).
\end{exercise}

\begin{exercise}\label{ex:3-2}
Derive the middle two lines of the rotated correlator, including the $\theta$ functions,
from the $\pm$ table. Then show directly from the trace identities that $D^{\qu\qu}=0$
and that $D^{\cl\qu}(t,s)$ vanishes for $s>t$, without using the table.
\end{exercise}

\begin{exercise}\label{ex:3-3}
Show that $S^\cl$ and $S^\qu$ of \cref{eq:SclSqu} obey the same two trace identities with
$\Tr_S$, and compute $\Tr_S[\tilde S^\cl(t)\tilde S^\cl(s)\rho_S]$ and
$\Tr_S[\tilde S^\cl(t)\tilde S^\qu(s)\rho_S]$ for $\rho_S=\tfrac12(\Id+r_z\sz)$ using
\cref{eq:S-int}. Interpret the second as the qubit's own response function.
\end{exercise}

\begin{exercise}\label{ex:3-4}
Verify \cref{eq:FDT} for the Lorentzian bath directly from \cref{eq:lorentz-kernels}, and
show that $\chi(\omega)=g^2\big[(\Omega-\omega-i\gamma/2)^{-1}+(\Omega+\omega+i\gamma/2)^{-1}\big]$.
\end{exercise}

\clearpage
